% Generated by scripts/materialize_unified_paper.py; do not edit.
\section{Controlled residual and strong concavity}\label{low:sec:residual}

Set
\[
 t_0=m_0\sqrt d,\qquad
 \gamma_d=\frac{m_0m'_0}{\sqrt d},\qquad
 e_{ij}=\theta(i-1).
\tag{L4.1}\label{low:main:eq:4-1}
\]
In decreasing order of the first endpoint define
\[
 T_{ij}=t_0+\gamma_d
 \sum_{\substack{q>i\\q\ne j}}T_{jq}+e_{ij}.
\tag{L4.2}\label{low:main:eq:4-2}
\]
The array is deterministic and adapted: when column \(i\) is exposed, every
future weight in \eqref{low:main:eq:4-2} has already been fixed.

With \(F_i\) as in \eqref{low:main:eq:3-1}, now using \(T\), we obtain the exact residual
identity
\[
 \partial_jF_i(c\mathbf1)=-\frac{T_{ij}-e_{ij}}{\sqrt d}.
\tag{L4.3}\label{low:main:eq:4-3}
\]
Thus the mismatch is explicit and deterministic.

\begin{lemma}[Uniform Hessian bound]\label{low:lem:hessian}
For all sufficiently large fixed \(C\), every
\(b\in[c-w,c+w]^{V_i}\) satisfies
\[
 -\nabla^2F_i(b)\succeq\mu_w I,
\tag{L4.4}\label{low:main:eq:4-4}
\]
where \(\mu_w>0\) is defined in \eqref{low:eq:mu-w}.
\end{lemma}

\begin{proof}
The recursion gives
\[
 T_{\max}\le\frac{t_0+\theta r}{1-\rho},\qquad
 \frac1d\max_j\sum_{q\ne j}T_{jq}
 \le\frac{m_0+\theta/D}{D(1-\rho)}.
\]
The diagonal of \(-\nabla^2F_i\) is at least \(j_-(w)\).  The absolute
off-diagonal row sum is at most the second term subtracted in
\eqref{low:eq:mu-w}, multiplied
by \(j_+(w)^2\).  Gershgorin's theorem proves \eqref{low:main:eq:4-4}.
\end{proof}

Write
\[
 \delta_j=b_j-c=\sqrt d\,h_{ij}+\varepsilon_{ij},\qquad
 |\varepsilon_{ij}|\le d^{-1/5}.
\]
Strong concavity, \eqref{low:main:eq:4-3}, and
\[
 \frac{e_{ij}\delta_j}{\sqrt d}-\frac{\mu_w\delta_j^2}{2}
 \le\frac{e_{ij}^2}{2\mu_wd}
\tag{L4.5}\label{low:main:eq:4-5}
\]
give
\[
F_i(b)\le F_i(c\mathbf1)-\sum_{j>i}T_{ij}h_{ij}
+\sum_{j>i}\frac{e_{ij}^2}{2\mu_wd}
+\frac{d^{-1/5}}{\sqrt d}\sum_{j>i}T_{ij}.
\tag{L4.6}\label{low:main:eq:4-6}
\]
The full boundary coefficient merges into the projection potential, while the
failure of exact gradient merge is paid by the square term.
